\chapter{S1：常温标定后独立验证}
\section{S1 测试项目总表}
\noindent\textbf{使用方法：} 先按本章进行测试计划评审，再按本章后续 CT-00--CT-18 和 CT-12A 的详细方案执行。所有测试在\emph{标定参数冻结、工况独立}前提下进行。
\small
\begin{longtable}{C{14mm}p{37mm}C{10mm}p{43mm}p{37mm}}
\toprule
编号 & 项目 & 优先 & 能力与测量对象 & 核心输出\\\midrule\endfirsthead
\toprule 编号 & 项目 & 优先 & 能力与测量对象 & 核心输出\\\midrule\endhead
CT-00 & 参考基准/装夹/数据链核查 & P0 & 静止/低速，所有实验前必做 & 基准、坐标、量程可用性\\
CT-01 & 参数生效与补偿前后对比 & P0 & 传感器六轴数据，参数只读 & 版本一致性、误差改善率\\
CT-02 & 常温零位与静态零速残差 & P0 & 夹具参考零位，ACC/GYRO & 水平角误差、静止 GYRO\\
CT-03 & ACC 三轴矢量独立验证 & P0 & 多个独立参考姿态 & 残差、RMSE、最大值\\
CT-04 & ACC 模长、对向与非线性诊断 & P0 & 正反姿态和中间角 & 模长误差、对向误差\\
CT-05 & 单方向静态倾角精度 & P0 & 已知参考静态姿态 & Roll/Pitch 各点误差\\
CT-06 & 复合姿态与交叉轴误差 & P0 & 内/外轴组合；非独立俯仰双轴 & 复合角度残差、耦合\\
CT-07 & 同姿态测量重复性 & P0 & 不卸装、多次到位 & 组内/组间标准差\\
CT-08 & 正反路径回程差 & P0 & 同一姿态从两方向接近 & 差值、最大回程差\\
CT-09 & GYRO 零速残差与姿态相关性 & P0 & 静止，多姿态 & 三轴残余零偏、方向依赖\\
CT-10 & GYRO 多速率正反向残差 & P0* & 速率转台；绝对精度需溯源 & 速率残差、正反对称性\\
CT-11 & GYRO 轴间耦合与低速线性 & P0* & 低速多方位转动 & 横向误差、回归残差\\
CT-12 & 常温静态噪声/短期漂移 & P0 & 固定姿态 30--60 分钟基线 & 标准差、漂移、均值波动\\
CT-12A & Allan 方差/随机稳定性表征 & P1 & 长时间静止连续采样（建议 3--12 小时） & ADEV、ARW、零偏不稳定性\\
CT-13 & 多次上电重复性 & P0 & 10 次（建议）独立电源循环 & 零偏/静态角复现性\\
CT-14 & 重复装夹与复测 & P1 & 卸装后重新定位（若允许） & 夹具误差、总复现性\\
CT-15 & 单次常温试验内温变化 & P1 & 20--30\,℃室温，记录内温 & 温度跨度与误差相关性\\
CT-16 & 低速运动后归零与恢复 & P1 & 安全路径运动-停止 & 回零差、恢复时间\\
CT-17 & 数据时间对齐与延迟诊断 & P1 & 小幅低速摆动、同步对比 & 延迟诊断及不确定性\\
CT-18 & 独立验证报告与追溯审核 & P0 & 结果/日志/证书复核 & 封存包与结论\\\bottomrule
\end{longtable}
\normalsize
\begin{keynote}
\textbf{P0* 特别说明：} CT-10、CT-11 对动态 IMU 内参非常重要，必须\emph{计划并实施}；但在现有 2TD-600 未取得速率参考的校准证书或等效溯源前，只能给出诊断结论，\textbf{不能出具陀螺仪绝对比例因子精度 PASS}。其余试验也需满足各自的参考不确定度要求。
\end{keynote}
